% Cote 119: offprints, reports, courses, CV and letters (typed and printed part): archive cote 119, site page 264 (PDF page 265).
% typescript. Draft compared: none.
% AI-transcribed by Claude Opus 5.5 (claude-opus-5-5, adaptive thinking, effort high), Message Batches API, batch r119_full (msgbatch_01DwZNm5Y7YsKBWyAN4sFs8w), 2026-11-13; not yet checked by eye.
\scanpage{119/264}
\mspage{11}

2. Le danger de sclérose d'un enseignant, qu'il soit simple assistant ou pro-%
fesseur, serait dans une large mesure évité par l'usage de l'année ``sabbatique''
pour les membres de l'Ecole de Mathématique, mentionné plus haut, et par
l'obligation faite à tout enseignant, de l'une ou l'autre section, d'un renou-%
vellement constant de son enseignement, et notamment des sujets enseignés,
tout au long de sa carrière. Il devrait être entendu notamment qu'un maître
de conférences ou professeur ne pourra enseigner le même cours plus de trois
années consécutives, et une \uline{commission d'arbitrage} à l'autorité incontestée
devrait pouvoir être saisie lorsque certains membres de l'Assemblée du dépar-%
tement jugeront que cette règle n'a pas été observée par un des enseignants.
Enfin, pour préserver un équilibre judicieux entre les deux sections du
département de Mathématique, nécessaire pour garantir la qualité de l'ensei-%
gnement en même temps que le niveau scientifique du département, on pourrait
prévoir une clause numérique, précisant que le nombre de maîtres de conférences
et de professeurs de chaque section doit être égal au tiers au moins du nombre
total des maîtres de conférences et professeurs.

3. La \uline{thèse de Mathématique} serait commune aux deux orientations ``recherche''
et ``enseignement''. Ce serait, comme par le passé, un travail écrit faisant
l'objet d'un exposé oral. Il devrait présenter certaines qualités d'originalité,
soit dans \unsure{\uline{la}|la} mise au point ou l'exposition d'une théorie mathématique assez
avancée, soit par la contribution de résultats originaux. Par le fait que des
contributions originales ne sont plus exigées (*), le niveau de la thèse propre-%
ment dite serait donc abaissé ; en revanche, la thèse serait un élément néces-%
saire mais non plus suffisant pour conférer le titre de docteur ès sciences
mathématiques. Pour celui-ci, on distingue deux ``mentions'', la \hand{\unsure{6}}mention \margin{\unsure{x}}
``recherche'' et la mention ``enseignement'', qui qualifient respectivement pour \margin{\unsure{$\lambda$}}
une nomination de maître de conférences à l'Institut de Mathématique et à
l'Ecole de Mathématique. Pour obtenir le doctorat mention recherche, il
faudra comme par le passé, en plus d'une culture mathématique solide, avoir
obtenu des résultats originaux d'une importance indiscutable. Compte tenu de
l'existence d'une carrière parallèle ``enseignement'', il conviendrait à ce
sujet d'être même plus exigeant qu'on ne l'a été souvent ces dernières années,

(*) Au lieu de supprimer entièrement cette exigence, on peut aussi essayer de
se borner à l'assouplir.

\margin{264}
